% ============================================================
%  文档类选项（在 main.tex 的 \documentclass[...]{tongjithesis} 中设置）
%  field=science    理工科（工科类、理科类，默认）— 章节编号：1 / 1.1 / 1.1.1
%  field=humanities 文科  （人文、法学、外语、艺术）        — 章节编号：一、/（一）/ 1.
%  经管类属社科：编号要求依理工，撰写依理工模版；26 届为过渡期，可选 science 或
%  humanities；自 27 届起统一使用 science。
% ============================================================

% ============================================================
%  封面信息
% ============================================================
\school{经济与管理学院}
\major{信息管理与信息系统}
\student{2253789}{张子越}
\thesistitle{考虑顾客选择的服务网络设计}{基于行程聚合与动态细化的精确优化方法}
\thesistitleeng{Service Network Design Considering Customer Choice}{An Exact Optimization Approach with Itinerary Aggregation and Dynamic Refinement}
\thesisadvisor{魏可佶\quad 副教授}
\thesisdate{\the\year}{\the\month}{\the\day} % 使用当前日期; 也可以自定义日期

% ============================================================
%  信息说明页
% ============================================================

% 成果类型（三选一）：
%   thesis      — 毕业论文，摘要
%   design      — 毕业设计（软件/创意/建筑等非工程设计类），摘要
%   engineering — 毕业设计（工程设计类），设计总说明
\infotype{thesis}

% 内容简述（请用300字以内简要概述）
\infoabstract{%
本文研究考虑顾客选择行为的服务网络设计问题，该问题由于候选行程集合规模爆炸并与其他变量耦合而难以求解。本文提出行程聚合与动态细化框架，该框架利用路径签名将候选行程集合聚合为紧凑的超级行程集合，在聚合网络上求解松弛模型获得下界，在完整网络上修复得到可行上界，并通过迭代细化保证有限收敛到全局最优。在真实航空排班实例上，超级行程聚合将142万条候选行程压缩至约21万条（减少85\%），框架在所有测试实例上均于5小时内达到低于1\%的最优性间隙，显著优于直接求解方法。
}

% 毕业设计：图纸数量与正文字数（成果类型为 design 或 engineering 时填写，两者须同时填写）
% \infodrawings{5}
% \infowordcount{15000}

% 毕业论文：正文总字数（成果类型为 thesis 时填写）
\infothesiswords{25681}

% 随附资料（每条用 \item 开头；不填则显示空白行）
\infomaterials{
  \item 附录
}
